% ---------------------------------------------------------------- 1 introduction
\section{Introduction}
\label{sec:intro}

The sixth generation (6G) of mobile networks is expected to support services whose value lies in the meaning of what is delivered rather than in the exact reproduction of every bit \cite{itu2023,strinati2021,letaief2019,saad2020}. Exact reproduction is the problem that Shannon solved \cite{shannon1948}; semantic communication addresses the level above it, identified by Weaver \cite{weaver1949}, by designing transmitter and receiver jointly around the information the receiver needs \cite{gunduz2023,luo2022,kountouris2021}. Learned joint source--channel coding has shown that such systems degrade gracefully with the channel and can outperform separate source and channel coding at short block lengths for text \cite{farsad2018,xie2021deepsc} and images \cite{bourtsoulatze2019,kurka2021}.

Large language models (LLMs) add a further resource: a strong model of what a sender is likely to say. When both ends share such a model, the receiver can combine the likelihood of the received symbols with the model's prediction of the next token and decide the most probable sentence, so that fewer channel symbols suffice for the same fidelity \cite{wang2024llmsc}; the context among tokens can likewise be exploited across modalities \cite{qiao2025token}. Pre-trained language models have also been used to rank the importance of words \cite{guo2023importance}, and large AI models are being studied as the knowledge base of semantic transceivers \cite{jiang2024largeai,xie2025largemodel,bariah2024}.

Two costs are usually left out of this picture. The first is the cost of keeping the shared knowledge aligned. After deployment the content changes---the topic of a conversation within minutes, the vocabulary of a domain over days---and a frozen receiver degrades. The receiver cannot adapt on its own, because only the transmitter knows the source text; every update of the receiver's model must therefore be carried over the very link the semantic codec is meant to economize. Existing work adapts or updates the knowledge of semantic transceivers after deployment \cite{zhang2023dynamic,zhang2026kbident,zhang2026kbupdate}; for a language-model prior updated at several timescales, however, the air time of the updates has not been set against the rate they save. The second is the cost of adapting the rate. Incremental redundancy lets a receiver ask for more symbols only when it needs them \cite{caire2001,jiang2022harq}, but every request travels on a return link, and a token-level policy makes many requests.

This article studies both costs in one system, called doubly nested semantic communication (DN-SC). The system is nested in rate: every token of the GPT-2 vocabulary \cite{radford2019} is mapped to a codeword of 16 real channel symbols whose leading coordinates are organized in shells, each normalized to unit power, so that any prefix is a valid transmission and sending a further shell after a prefix is exactly equivalent to having sent the longer prefix at once. The receiver decodes with the language model as its prior and requests further shells only for the tokens whose posterior confidence is below a threshold. The system is also nested in time: the receiver's prior is shaped by three memories that change at three speeds---the context of the conversation (every message, free of signaling), a logit bias in embedding space (every ten adaptation steps) and low-rank adapter slots of the language model \cite{hu2022lora} (every fifty steps, with one slot per knowledge domain). The multi-frequency organization follows the nested-learning view of models as systems of components updated at different rates \cite{behrouz2025}; in a transceiver, however, each update frequency acquires a price in channel uses, and DN-SC records that price exactly.

The evaluation is designed so that the results can contradict the method. Teacher-forced emulation, which scores each token given the true tokens before it, is cheap but optimistic; on a subset of steps the same received symbols are therefore decoded for real, free-running, and a decision rule fixed in advance determines which measurement stands. The decoder of the refinement policy was chosen on validation text only, by a rule written down before it was run, and the confirmatory comparison uses six seeds, paired bootstrap intervals and an exact seed-level test.

The contributions are as follows.
\begin{enumerate}
\item \textbf{A rate-nested token codebook with exact power control.} Each shell of each message is scaled to unit power, so every prefix meets the power constraint and token-level incremental redundancy is exact. Nesting costs nothing measurable at the prefix that was checked: at 8 symbols per token the nested codebook reaches a channel-only token accuracy of 0.8885 against 0.8890 for a codebook trained for 8 symbols alone.
\item \textbf{A time-nested receiver with an exact synchronization ledger.} The context memory is built by the receiver from its own decisions and costs nothing on the air; the logit-bias and adapter memories are learned at the transmitter from the true text and shipped as quantized differences at their own periods. Adapter slots keep one adapter per domain, and every shipped bit is counted.
\item \textbf{A complete account in channel symbols.} Forward symbols, the return-link signaling of the refinement requests (as an explicit bitmap and at its entropy) and the synchronization of the memories are all converted to real channel symbols, giving a net semantic rate and the traffic volume above which shipping adaptation pays.
\item \textbf{A pre-specified evaluation and its result.} On a hierarchical drift stream with six seeds, DN-SC needs 23.3\,\% fewer real channel symbols per token than a frozen fixed-rate receiver with the same decoder at 90\,\% token accuracy, with the request signaling included, and 27.9\,\% fewer when the audited messages are decoded for real (the synchronization of the memories is accounted separately); the decomposition shows that the free context memory carries most of the gain of the memories and that shipped adaptation pays only above a large traffic volume.
\end{enumerate}
The code, the configurations and every result file behind the tables are public (see the Data Availability Statement).

The remainder of the article is organized as follows. Section~\ref{sec:related} reviews related work. Section~\ref{sec:system} states the system model and the problem. Section~\ref{sec:method} presents the DN-SC receiver, its memories, the refinement protocol and the cost accounting. Section~\ref{sec:setup} describes the experimental protocol, Section~\ref{sec:results} reports and discusses the results, and Section~\ref{sec:conclusion} concludes.
